% ch11_limitations.tex: Chapter 11, Limitations and threats to validity.
% Collects the limits stated in Chapters 4 to 10, one sentence each. Every number is traced in ch11_limitations_NOTES.md.
\chapter{Limitations and threats to validity}\label{ch:limitations}

This chapter collects the limits stated in the earlier chapters, one sentence each, so that every claim of Table~\ref{tab:T20} can
be weighed against them.

\section{The sealed E1 result}\label{sec:lim-e1}

\begin{itemize}
  \item The one confirmed claim rests on one sealed set and one primary contrast, with a lower end of +0.0013 and a derived smallest
  detectable effect of about 0.074, so E1 supports a gain without pinning down its size.
  \item The claim holds for the recipe averaged over seeds 5 to 9, not for one checkpoint: seed 5 alone is not separated from C1
  (+0.0189 [$-$0.0403, +0.1068]).
  \item E1 was read at OP-A only, a stated deviation from the protocol, with width guards in place of OP-W.
  \item 101 of the 106 E1 cells are CSIR, and the 5 BMRK cells support no statement.
  \item The full cells carry the gain; on the 62 half-cut cells it is not detected (+0.0110 [$-$0.0289, +0.0722]).
  \item ``Unseen modules'' is supported, not proven: only about a third of the image-id pairs could be tested, and half-cut images
  and BMRK cells had no testable counterpart.
  \item The E1 cells had been looked at before (7 registered openings, 3 of them at one contrast, reads before 31 October and one
  use of published E1 contrasts), although the final models had not.
  \item E1 separates the model from the D4 views but not HRNet from SWAD, and the D4 part is not resolved on E1.
  \item Inactive is not decided on E1 (+0.0996 [$-$0.0310, +0.2162], 10 inactive cells).
  \item Both arms predict under half of the labelled E1 crack pixels (width ratios 0.458 and 0.463), and the read cannot tell wider
  E1 labels from missed cracks.
  \item After opening \#9, E1 is training data for the few-shot models, so a new unseen-module claim needs a new labelled set.
\end{itemize}

\section{The CV5 evidence}\label{sec:lim-cv5}

\begin{itemize}
  \item The \texttt{swad} readout and D4 were chosen on pool cells, so every CV5 number of the C2 recipe is selection-biased and its
  CV5 pass (+0.0222 [+0.0138, +0.0308]) is no confirmation.
  \item The full-pool models C1 and C2 have no held-out score of their own.
  \item No HOUSE prediction with D4 exists, so the HOUSE rule sets the eight-view C2 recipe against a single-view HOUSE, and the
  single-view CS $-$ HOUSE holds 0 at OP-A (+0.0107 [$-$0.0009, +0.0229]).
  \item C2 costs $11.8\times$ C1 in FLOPs and no equal-cost comparator was printed for it; the same-view controls are not equal cost.
  \item The margin $M$ covers seed noise, not the choice of folds; the split-variance probe was never run.
  \item The passing crack contrasts rest on few module groups: removing 17 (HRNet) to 31 (HRNet + SWAD) of the 125 groups with
  crack pixels takes them to 0.
  \item The 2-seed screens ran on cv5v2, which leaked modules across folds, and on two machines, so they form a separate tier read
  per machine.
\end{itemize}

\section{Labels and metrics}\label{sec:lim-labels}

\begin{itemize}
  \item Every pool cell has a single annotation; the one double annotation (strict crack IoU 0.259, relF1@2 0.829) suggests, but
  does not measure, how far labels limit any score.
  \item Crack label errors are a minority of the image-supported candidates (corrected share 0.244 [0.105, 0.436]), but the
  detector's recall is unknown and the raters were runs of one model, not independent people.
  \item Strict crack IoU measures the line-width convention (PF 0.233, 0.365 and 0.414 under three conventions), so no strict-IoU
  number here decides anything.
  \item Label width is a source convention (ARTS 1.59 to BMRK 3.09~px per centreline pixel), so the per-source width guard partly
  measures the annotators.
  \item E1 labels are about $1.5\times$ wider than the pool labels, so pool-trained models meet a different convention on E1.
  \item The 2~px tolerance of relF1@2 rests on the few duplicate annotations available, and relF1@2 is blind to line width without
  its guard.
  \item The images are windows of about $3\times3$ cells (57\% of crack and 71\% of inactive pixels lie outside the centre cell),
  so every cell-level quantity describes a window.
  \item Inactive pixels sit in 33 of the 285 module groups (Kish effective $n$ 12.9), so inactive contrasts carry intervals of about
  $\pm$0.05 to $\pm$0.06 (audit estimate).
\end{itemize}

\section{Data, test50 and licences}\label{sec:lim-data}

\begin{itemize}
  \item test50 is in-distribution with known leakage (8 cells are the same physical cell as a pool cell, 7 by pixels and 1 by the
  cell-table link only; 7 are neighbouring crops; 19 share a module) and had 19 counted openings, so its numbers are context only
  and support no claim to beat a published model.
  \item BMRK and SDLE cells cannot be grouped by physical module, so leakage between their neighbouring windows cannot be excluded.
  \item All labelled cells come from one dataset family, and the source is recognisable from global image statistics (balanced
  accuracy 0.94), a shortcut open to any model.
  \item The weak classes are not solved: material and edge dark score 0.0000 on E1 for every arm, and OP-C, which raised mIoU-R13
  on the pool (+0.0126 [+0.0077, +0.0167]), is in no final model.
  \item The 50 BMRK cells are ELPV images under CC BY-NC-SA 4.0 and BenchmarkEL has no licence file, so C2 is research-only as
  trained, and retraining without the BMRK cells was not measured.
  \item No WattCycle module was imaged, so the model is trained and tested on public data only.
\end{itemize}

\section{Cost, severity and the few-shot track}\label{sec:lim-other}

\begin{itemize}
  \item All latencies were measured on shared hosts and cover network time only, so they are upper bounds for idle hardware.
  \item The severity output was validated only for PF out of fold on pool cells, not for C2, on E1 or on half-cut cells.
  \item Predicted inactive area is not shown to predict power loss better than crack length, and no watt figure is given.
  \item The few-shot track rests on about 9 effective images, has 80\% power only for gains of about +0.085 or more, and may be
  optimistic for a new site because its folds group by source image; its confirmed gain (+0.2213 [+0.0977, +0.3504]) is a
  few-shot result inside E1 (3 seeds, 101 of 106 cells CSIR), says nothing about a new layout without labels, and with ALL $-$ HALF
  not separated (+0.0171 [$-$0.0246, +0.0625]) does not say how many labelled cells a layout needs.
  \item The few-shot models learn the about $1.5\times$ wider E1 labels (width ratio ALL / ZS 1.426, reported), which lifts strict
  IoU and pixel recall without finding more cracks; relF1@2 decides because it is less sensitive to this, though not immune.
\end{itemize}
